% Irodov Part VI §6.2, original 66--80. Printed pages 239--240.
% User-provided Arihant PDF pages 248--249; reference answers PDF p.350.
% Uncertainty estimates explicitly distinguish hbar and Delta x conventions.
\wavequantproblem{66}{239}
已知电子、质子以及质量为 $m_b=1\,\mathrm{mg}$ 的小球，其位置（小球为质心位置）的不确定度均为 $\Delta x=1\,\mu\mathrm m$。估算测量其速度时的最小不确定度。
\exsolution
海森堡不确定关系
$\Delta x\,\Delta p_x\geq\hbar/2$，
对低速物体 $\Delta p_x=m\Delta v_x$，故
\[
 \Delta v_{x,\min}=\frac{\hbar}{2m\Delta x}.
\]
取 $m_e=9.109\times10^{-31}\,\mathrm{kg}$、
$m_p=1.673\times10^{-27}\,\mathrm{kg}$、
$m_b=10^{-6}\,\mathrm{kg}$，得
\[
 \boxed{\Delta v_{e,\min}=5.79\times10^1\,\mathrm{m/s}},
 \quad\boxed{\Delta v_{p,\min}=3.15\times10^{-2}\,\mathrm{m/s}},
 \quad\boxed{\Delta v_{b,\min}=5.27\times10^{-23}\,\mathrm{m/s}}.
\]
\excheck 若原书按数量级用 $\Delta x\Delta p\sim\hbar$，
数值可差一个2的系数；三者的速度精度差异随质量倒数变化，
这不影响微观与宏观比较结论。

\wavequantproblem{67}{239}
利用不确定原理估算电子在氢原子中的速度不确定度，取原子线尺度 $l=0.10\,\mathrm{nm}$，并与第一玻尔轨道的电子速度比较。
\exsolution
当电子的位置被限制在 $l$ 的量级，
动量不确定度在数量级上为
$\Delta p\sim\hbar/l$，所以
\[
 \boxed{\Delta v\sim\frac{\hbar}{m_e l}
 \simeq1.16\times10^6\,\mathrm{m/s}}.
\]
第一玻尔轨道电子速度
$v_1=\alpha_{\rm fs}c\simeq2.19\times10^6\,\mathrm{m/s}$，
从而
\[
 \boxed{\frac{\Delta v}{v_1}\sim0.53}.
\]
\excheck 本题为数量级估算，而非通过精确几率分布求标准差；
电子速度的不确定量确实与其典型轨道速度同阶。

\wavequantproblem{68}{239}
证明当粒子的坐标不确定度为 $\Delta x=\lambda/(2\pi)$，
其中 $\lambda$ 为其德布罗意波长时，
它的速度不确定度与粒子速度本身同数量级。
\exsolution
由 $\lambda=h/p=2\pi\hbar/p$ 得
$\Delta x=\hbar/p$。
不确定性关系给出
\[
 \Delta p\geq\frac{\hbar}{2\Delta x}=\frac p2.
\]
若粒子处于非相对论区，
$p=mv$、$\Delta p=m\Delta v$，
于是
\[
 \boxed{\Delta v\geq\frac v2}.
\]
\excheck 系数 $1/2$ 依具体标准差定义；
结论是$\Delta v$不能比$v$小很多，即为同阶。

\wavequantproblem{69}{240}
一个自由电子最初局限在长度量级
$l=0.10\,\mathrm{nm}$ 的空间区域，
用不确定原理估计其波包宽度增至原来
$\eta=10$ 倍所需时间。
\exsolution
初始局域化使动量展宽在数量级上
$\Delta p\sim\hbar/l$，
速度展宽为 $\Delta v\sim\hbar/(m_e l)$。
自由传播中不同速度的波包分量逐渐分离，
波包额外扩展距离约 $\Delta v\,t$。
当展宽达到 $\eta l$ 的量级时
\[
 \eta l\sim\frac{\hbar t}{m_e l},
 \qquad
 \boxed{t\sim\frac{\eta m_e l^2}{\hbar}
 \simeq8.6\times10^{-16}\,\mathrm s}.
\]
\excheck 这是数量级估算，精确高斯波包的展开时间还取决于
$l$ 是振幅宽度、概率标准差还是其他宽度约定；
不宜将题中的“估计”误写成唯一精确时间。

\wavequantproblem{70}{240}
利用不确定原理估算局域于
$l=0.20\,\mathrm{nm}$ 的电子的最小动能。
\exsolution
电子被束缚在长度$l$量级时，
$\Delta p\sim\hbar/l$，
取典型动量$p\sim\Delta p$，
\[
 \boxed{T_{\min}\sim\frac{\hbar^2}{2m_e l^2}
 \simeq0.95\,\mathrm{eV}\approx1\,\mathrm{eV}}.
\]
\excheck 上式是原书要求的粗估下限量级；
若指定精确无限深一维方阱边界，
其基态能级则为$\pi^2\hbar^2/(2m_el^2)$，
有额外的$\pi^2$系数。

\wavequantproblem{71}{240}
电子动能约 $T=4\,\mathrm{eV}$，
被局限在尺寸$l=1\,\mu\mathrm m$的区域。
估算相对速度不确定度。
\exsolution
由不确定关系按量级
$\Delta p\sim\hbar/l$；
非相对论动量$p=\sqrt{2m_eT}$，
所以
\[
 \boxed{\frac{\Delta v}{v}\sim\frac{\Delta p}{p}
 \sim\frac{\hbar}{l\sqrt{2m_eT}}
 \simeq9.8\times10^{-5}\approx10^{-4}}.
\]
\excheck 由于$l$远大于电子的德布罗意波长，
相对速度展宽远小于1，量级合理。

\wavequantproblem{72}{240}
电子位于宽度为$l$的无限深一维方阱，
用不确定原理估算占据最低可能能量时
电子施加在阱壁上的力。
\exsolution
最小动量量级为
$p\sim\hbar/l$，
相应能量为
$E_{\min}\sim\hbar^2/(2m_el^2)$。
缓慢改变壁间距所需的功
$dW=F\,dl$，
故电子对可移动阱壁的压力效应
\[
 \boxed{F\sim-\frac{dE_{\min}}{dl}
 \sim\frac{\hbar^2}{m_el^3}}.
\]
\excheck 精确无限阱基态是
$E_1=\pi^2\hbar^2/(2m_el^2)$，
对应的精确广义力为
$F_1=\pi^2\hbar^2/(m_el^3)$；
原书要求的是不确定原理估算。

\wavequantproblem{73}{240}
质量为$m$的粒子在一维简谐势
$U(x)=kx^2/2$中运动。
用不确定原理估计最低允许能量，
并与精确解比较。
\exsolution
若位置展宽为$x$的量级，
则$p\sim\hbar/x$，因而估计总能量
\[
 E(x)\sim\frac{\hbar^2}{2mx^2}+\frac{kx^2}{2}.
\]
极小值条件
$dE/dx=-\hbar^2/(mx^3)+kx=0$，
给出
$x^4=\hbar^2/(mk)$，
于是
\[
 \boxed{E_{\min}\sim\hbar\sqrt{k/m}
 =\hbar\omega},\qquad \omega=\sqrt{k/m}.
\]
\excheck 采用真正的标准差不确定关系与高斯态变分
可以得到精确基态能
$E_0=\hbar\omega/2$；旧式量级估算只保证尺度。

\wavequantproblem{74}{240}
用不确定原理估计氢原子电子所允许的
最低能量及对应的最概然量级核间距离。
\exsolution
电子局域到尺度$r$，按量级
$p\sim\hbar/r$；总能量近似
\[
 E(r)\sim\frac{\hbar^2}{2m_er^2}
 -\frac{k_ee^2}{r},
 \qquad k_e=\frac1{4\pi\varepsilon_0}.
\]
令$dE/dr=0$，
$-\hbar^2/(m_er^3)+k_ee^2/r^2=0$，得
\[
 \boxed{r_*\sim\frac{\hbar^2}{m_ek_ee^2}
 =a_0\simeq52.9\,\mathrm{pm}}.
\]
代回能量
\[
 \boxed{E_{\min}\sim-\frac{m_e(k_ee^2)^2}{2\hbar^2}
 \simeq-13.6\,\mathrm{eV}}.
\]
\excheck 用等式求极小值只是在估算有效尺度；
真正三维氢基态径向概率最大点恰为$a_0$，
但这需用量子波函数另行证明。

\wavequantproblem{75}{240}
速度$v=600\,\mathrm{m/s}$的平行氢原子束
垂直照射狭缝，缝后距离$\ell=1.0\,\mathrm m$
处放置屏。用不确定原理求成像总宽度最小时
狭缝宽度$\delta$。
\exsolution
狭缝横向局域化给
$\Delta p_y\sim\hbar/\delta$。
原子纵向动量$p_x\simeq m_Hv$，
角展宽$\Delta\theta\sim\Delta p_y/p_x$，
传播距离$\ell$引起的附加宽度为
$\Delta_{\rm dif}\sim\ell\hbar/(m_Hv\delta)$。
总宽度估算
\[
 W(\delta)\sim\delta+\frac{\hbar\ell}{m_Hv\,\delta}.
\]
对$\delta$求极值，得
\[
 \boxed{\delta_{\min}\sim\sqrt{\frac{\hbar\ell}{m_Hv}}
 \simeq1.03\times10^{-5}\,\mathrm m
 \approx0.01\,\mathrm{mm}}.
\]
\excheck 增大缝宽会直接增加几何模糊，
减小缝宽则会增加衍射模糊，因此存在非零最优值。

\wavequantproblem{76}{240}
求一质量为$m$的自由粒子的
含时薛定谔方程的一个特解。
\exsolution
自由粒子$U=0$，满足
\[
 i\hbar\frac{\partial\Psi}{\partial t}
 =-\frac{\hbar^2}{2m}\nabla^2\Psi.
\]
取平面波试探解
$\Psi=A\exp[i(\boldsymbol k\cdot\boldsymbol r-\omega t)]$。
代入得
$\hbar\omega=\hbar^2k^2/(2m)$，
所以一个特解是
\[
 \boxed{\Psi(\boldsymbol r,t)=
 A\exp\!\left[\frac{i}{\hbar}
 \left(\boldsymbol p\cdot\boldsymbol r-\frac{p^2}{2m}t\right)\right]},
 \quad\boldsymbol p=\hbar\boldsymbol k.
\]
\excheck 无限空间中的单个平面波不是平方可积函数，
描述可归一化波包时须作动量态的叠加。

\wavequantproblem{77}{240}
宽度为$l$的无限深一维势阱$0<x<l$
内粒子处于基态。求它位于
$l/3\le x\le2l/3$的概率。
\exsolution
归一化基态波函数
$\psi_1(x)=\sqrt{2/l}\sin(\pi x/l)$。
因此
\[
 P=\frac2l\int_{l/3}^{2l/3}\sin^2\frac{\pi x}{l}\,dx.
\]
利用$\sin^2u=(1-\cos2u)/2$，
\[
 \boxed{P=\frac13+\frac{\sqrt3}{2\pi}
 \simeq0.609}.
\]
\excheck 概率超过几何区间比例$1/3$，
因为基态概率密度在阱中部最大。

\wavequantproblem{78}{240}
无限深一维方势阱以阱中心为坐标原点，
宽度为$l$，即$|x|<l/2$。
求各定态的归一化波函数。
\exsolution
阱内定态方程
$\psi''+k^2\psi=0$；
边界$\psi(\pm l/2)=0$，
故$k=n\pi/l$、$n=1,2,\dots$。
按奇偶性写成
\[
 \boxed{\psi_n(x)=
 \begin{cases}
 \sqrt{2/l}\cos(n\pi x/l),&n=1,3,5,\ldots,\\
 \sqrt{2/l}\sin(n\pi x/l),&n=2,4,6,\ldots,
 \end{cases}}
\quad |x|<l/2,
\]
阱外$\psi_n=0$。
对应能量
$E_n=n^2\pi^2\hbar^2/(2ml^2)$。
\excheck 余弦/正弦只差可忽略的整体符号；
归一化需满足$\int_{-l/2}^{l/2}|\psi_n|^2dx=1$。

\wavequantproblem{79}{240}
证明无限深一维方阱$0<x<l$
内不同能级的定态波函数彼此正交。
\exsolution
取标准归一化本征函数
$\psi_n(x)=\sqrt{2/l}\sin(n\pi x/l)$，
$n=1,2,\dots$。
对$n\ne n'$，用积化和差：
\[
 \int_0^l\psi_n\psi_{n'}dx
 =\frac1l\int_0^l
 \left[\cos\frac{(n-n')\pi x}{l}
 -\cos\frac{(n+n')\pi x}{l}\right]dx=0.
\]
因此
\[
 \boxed{\langle n|n'\rangle=\delta_{nn'}}.
\]
\excheck $n=n'$时积分为1；
正交性可独立由厄米哈密顿算符不同本征值证明。

\wavequantproblem{80}{240}
电子位于长$l$的无限深一维势阱，
能级极为密集。推导单位能量间隔内能级数
$dN/dE$，并求$l=1.0\,\mathrm{cm}$、
$E=1.0\,\mathrm{eV}$时的值。
\exsolution
精确能级
$E_n=n^2\pi^2\hbar^2/(2m_el^2)$。
高$n$连续近似下
$n(E)=l\sqrt{2m_eE}/(\pi\hbar)$。
对$E$求导，得
\[
 \boxed{\frac{dN}{dE}=\frac l{\pi\hbar}
 \sqrt{\frac{m_e}{2E}}}.
\]
此式以J为能量单位。若改为每eV的态数，
须乘$1\,\mathrm{eV}=1.602176634\times10^{-19}\,\mathrm J$，
故
\[
 \boxed{\left.\frac{dN}{dE}\right|_{1\,\mathrm{eV}}
 \simeq8.1\times10^6\ \mathrm{eV^{-1}}}.
\]
\excheck 连续密度只适用于$E_n$相邻间隔远小于所考察能量范围，
且这里未计电子自旋二重简并。
